
\section*{Chapitre 8 - Notion de fraction}
\addcontentsline{toc}{section}{Chapitre 8 - Notion de fraction}

\subsection*{I. Fractions et partage}
\addcontentsline{toc}{subsection}{I. Fractions et partage}

\begin{Définitions}
Pour tout entier $b \neq 0$, la \textbf{fraction} $\dfrac{1}{b}$ correspond à une part de l'unité lorsque l'on partage cette unité en $b$ parts égales. Pour tout entier $a$, $\dfrac{a}{b} = \dfrac{1}{b} + \dfrac{1}{b} + ... + \dfrac{1}{b} = a \times \dfrac{1}{b}$. \\

Le nombre $a$ est appelé le \textbf{numérateur} de la fraction et le nombre $b$ est le \textbf{dénominateur}.
\end{Définitions}

\begin{Remarque}
Si $b \neq 0$, on voit donc que $b \times \dfrac{1}{b} = \dfrac{b}{b} = 1$.
\end{Remarque}


\begin{Exemples}
Quelle fraction de la figure est colorée ?
\vspace{-0.5cm}
\begin{multicols}{3}
\begin{itemize}[itemsep=1em]
	\item $\ldots$ :
	\begin{minipage}[c]{\linewidth} \begin{tikzpicture}[baseline,scale = 0.5]

    \tikzset{
      point/.style={
        thick,
        draw,
        cross out,
        inner sep=0pt,
        minimum width=5pt,
        minimum height=5pt,
      },
    }
    \clip (-0.5,-0.5) rectangle (5,3.5);
    	\draw[color={black}] (0,0)--(4.5,0)--(4.5,3)--(0,3)--cycle;
	\draw[color={black},line width = 0,preaction={fill,color = {Couleur10}, opacity = 0.3}] (0,0)--(0,3)--(1.125,3)--cycle;
	\draw[color={black},line width = 0,preaction={fill,color = {Couleur10}, opacity = 0.3}] (0,0)--(1.125,0)--(1.125,3)--cycle;
	\draw[color={black},line width = 0,preaction={fill,color = {Couleur10}, opacity = 0.3}] (2.25,0)--(2.25,3)--(3.375,3)--cycle;
	\draw[color ={black}] (0,0)--(1.13,3);
	\draw[color ={black}] (0,3)--(1.13,3);
	\draw[color ={black}] (0,0)--(1.13,3);
	\draw[color ={black}] (1.13,0)--(1.13,3);
	\draw[color ={black}] (1.13,0)--(2.25,3);
	\draw[color ={black}] (1.13,3)--(2.25,3);
	\draw[color ={black}] (1.13,0)--(2.25,3);
	\draw[color ={black}] (2.25,0)--(2.25,3);
	\draw[color ={black}] (2.25,0)--(3.38,3);
	\draw[color ={black}] (2.25,3)--(3.38,3);
	\draw[color ={black}] (2.25,0)--(3.38,3);
	\draw[color ={black}] (3.38,0)--(3.38,3);
	\draw[color ={black}] (3.38,0)--(4.5,3);
	\draw[color ={black}] (3.38,3)--(4.5,3);
	\draw[color ={black}] (3.38,0)--(4.5,3);
	\draw[color ={black}] (4.5,0)--(4.5,3);

\end{tikzpicture}
 \end{minipage}
	\item $\ldots$ :
	\begin{minipage}[c]{\linewidth} \begin{tikzpicture}[baseline,scale = 0.5]

    \tikzset{
      point/.style={
        thick,
        draw,
        cross out,
        inner sep=0pt,
        minimum width=5pt,
        minimum height=5pt,
      },
    }
    \clip (-2.4,-2.12) rectangle (2.4021130325903073,2.5);
    	\draw[color ={black}] (-1.18,-1.62)--(1.18,-1.62);
	\draw[color={black},line width = 0,preaction={fill,color = {Couleur10}, opacity = 0.3}] (0,0)--(1.175571,-1.618034)--(1.902113,0.618034)--cycle;
	\draw[color ={black}] (1.18,-1.62)--(1.9,0.62);
	\draw[color ={black}] (1.9,0.62)--(0,2);
	\draw[color ={black}] (0,2)--(-1.9,0.62);
	\draw[color ={black}] (-1.9,0.62)--(-1.18,-1.62);
	\draw[color ={black}] (0,0)--(-1.18,-1.62);
	\draw[color ={black}] (0,0)--(1.18,-1.62);
	\draw[color ={black}] (0,0)--(1.9,0.62);
	\draw[color ={black}] (0,0)--(0,2);
	\draw[color ={black}] (0,0)--(-1.9,0.62);

\end{tikzpicture}
 \end{minipage}
	\item $\ldots$ :
	\begin{minipage}[c]{\linewidth} \begin{tikzpicture}[baseline,scale = 0.5]

    \tikzset{
      point/.style={
        thick,
        draw,
        cross out,
        inner sep=0pt,
        minimum width=5pt,
        minimum height=5pt,
      },
    }
    \clip (-2.6,-2.6) rectangle (2.5,2.6);
    	\draw  [color={black},line width = 0,preaction={fill,color = {Couleur10},opacity = 0.3}] (4.440892098500626e-16,2) -- (0,0) -- (1.4142135623730951,1.414213562373095) arc (45:90:2) ;
	\draw  [color={black},line width = 0,preaction={fill,color = {Couleur10},opacity = 0.3}] (-1.414213562373095,1.4142135623730951) -- (0,0) -- (1.2246467991473532e-16,2) arc (90:135:2) ;
	\draw  [color={black},line width = 0,preaction={fill,color = {Couleur10},opacity = 0.3}] (-2,4.440892098500626e-16) -- (0,0) -- (-1.414213562373095,1.4142135623730951) arc (135:180:2) ;
	\draw  [color={black},line width = 0,preaction={fill,color = {Couleur10},opacity = 0.3}] (-1.4142135623730954,-1.4142135623730947) -- (0,0) -- (-2,2.4492935982947064e-16) arc (180:225:2) ;
	\draw  [color={black},line width = 0,preaction={fill,color = {Couleur10},opacity = 0.3}] (-4.440892098500626e-16,-2) -- (0,0) -- (-1.4142135623730954,-1.414213562373095) arc (-135:-90:2) ;
	\draw  [color={black},line width = 0,preaction={fill,color = {Couleur10},opacity = 0.3}] (1.4142135623730947,-1.4142135623730954) -- (0,0) -- (-3.6739403974420594e-16,-2) arc (-90:-45:2) ;
	\draw[color={black},fill opacity = 1.1] (0,0) circle (2);
	\draw[color ={black}] (0,0)--(2,0);
	\draw[color ={black}] (0,0)--(1.41,1.41);
	\draw[color ={black}] (0,0)--(0,2);
	\draw[color ={black}] (0,0)--(-1.41,1.41);
	\draw[color ={black}] (0,0)--(-2,0);
	\draw[color ={black}] (0,0)--(-1.41,-1.41);
	\draw[color ={black}] (0,0)--(0,-2);
	\draw[color ={black}] (0,0)--(1.41,-1.41);

\end{tikzpicture}
 \end{minipage}
\end{itemize}
\end{multicols}


\end{Exemples}

\begin{Propriété}
La fraction $\dfrac{a}{b}$, avec $b \neq 0$, est le nombre qui, multiplié par $b$, donne $a$ : $\dfrac{a}{b} \times b = b \times \dfrac{a}{b} = a$.
\end{Propriété}


\begin{Exemple}
La fraction $\dfrac{2}{3}$ est le nombre qui, multiplié par 3, donne 2 ; c'est-à-dire $\dfrac{2}{3} \times 3 = 3 \times \dfrac{2}{3} = 2$. \\ 
$ $ \\
\begin{minipage}[c]{0.8\textwidth}
\begin{center}
\begin{tikzpicture}[scale=0.9]
    % Définir les hauteurs des rectangles
    \def\tallrowheight{1.2}       % Pour les lignes avec fractions
    \def\shortrowheight{0.8}      % Pour la première ligne (plus courte)
    \def\totalheight{2*\tallrowheight + \shortrowheight}
    
    % Premier rang du tableau (tons de violet)
    \fill[Couleur8!20] (0,\totalheight-\shortrowheight) rectangle (12,\totalheight);
    \draw[Couleur8,very thick] (0,\totalheight-\shortrowheight) rectangle (12,\totalheight);
    \draw[Couleur8,very thick] (6,\totalheight-\shortrowheight) -- (6,\totalheight);
    
    % Deuxième rang du tableau (tons de bleu)
    \fill[Couleur6!20] (0,\tallrowheight) rectangle (12,\totalheight-\shortrowheight);
    \draw[Couleur8,very thick] (0,\tallrowheight) rectangle (12,\totalheight-\shortrowheight);
    \foreach \x in {2,4,6,8,10}
        \draw[Couleur8,very thick] (\x,\tallrowheight) -- (\x,\totalheight-\shortrowheight);
    
    % Troisième rang du tableau (tons de vert)
    \fill[Couleur7!15] (0,0) rectangle (12,\tallrowheight);
    \draw[Couleur8,very thick] (0,0) rectangle (12,\tallrowheight);
    \draw[Couleur8,very thick] (4,0) -- (4,\tallrowheight);
    \draw[Couleur8,very thick] (8,0) -- (8,\tallrowheight);
    
    % Lignes du tableau (ajustées pour couvrir toute la largeur)
    \draw[Couleur8,very thick] (0,0) -- (12,0);
    \draw[Couleur8,very thick] (0,\tallrowheight) -- (12,\tallrowheight);
    \draw[Couleur8,very thick] (0,\totalheight-\shortrowheight) -- (12,\totalheight-\shortrowheight);
    \draw[Couleur8,very thick] (0,\totalheight) -- (12,\totalheight);
    
    % Contenu du premier rang (en violet foncé)
    \foreach \x in {3,9}
        \node[text=Couleur8] at (\x, \totalheight-\shortrowheight/2) {1};
    
    % Contenu du deuxième rang (en bleu)
    \foreach \x in {1, 3, 5, 7, 9, 11}
        \node[text=Couleur6] at (\x, \totalheight-\shortrowheight-\tallrowheight/2) {$\dfrac{1}{3}$};
    
    % Contenu du troisième rang (en vert)
    \node[text=Couleur7] at (2, \tallrowheight/2) {$\dfrac{2}{3} = 2 \times \dfrac{1}{3}$};
    \node[text=Couleur7] at (6, \tallrowheight/2) {$\dfrac{2}{3}$};
    \node[text=Couleur7] at (10, \tallrowheight/2) {$\dfrac{2}{3}$};
\end{tikzpicture}
\end{center}
\end{minipage} \hfill
\begin{minipage}[c]{0.2\textwidth}
\begin{center}
1 + 1 = 2 \\
$ $ \\
$ $ \\
$ $ \\
$ \dfrac{2}{3} \times 3 = 2 $
\end{center}
\end{minipage}

\end{Exemple}

\subsection*{II. Fractions et quotients}
\addcontentsline{toc}{subsection}{II. Fractions et quotients}


\begin{Propriété}
Pour deux nombres entiers $a$ et $b$ avec $b \neq 0$, la fraction $\dfrac{a}{b}$ est le \textbf{résultat exact} de la division de $a$ par $b$. C'est le \textbf{quotient} de $a$ par $b$, autrement dit, $\dfrac{a}{b} = a \div b$.
\end{Propriété}

\begin{Exemples}
Donner l'écriture décimale des quotients : 
\vspace{-0.5cm}
\begin{multicols}{4}
\begin{itemize}
\item $\dfrac{3}{10} = \ldots$
\item $\dfrac{1}{2} = \ldots$
\item $\dfrac{3}{5} = \ldots$
\item $\dfrac{1}{4} = \ldots$
\end{itemize}
\end{multicols}
\end{Exemples}

\subsection*{III. Fractions et demi-droites graduées}
\addcontentsline{toc}{subsection}{III. Fractions et demi-droites graduées}

\begin{Exemple}
\begin{enumerate}
\item Placer le point $A$ d'abscisse $\dfrac{7}{3}$, le point $B$ d'abscisse $\dfrac{1}{3}$ et le point $C$ d'abscisse $\dfrac{2}{3}$.


\bigskip
  \begin{tikzpicture}[x=2.5mm]
  \draw[-{Latex[round]},thick] (0,0) -- (61,0);
  \foreach \x in {0,3.33333,...,60} \draw[thick] ([yshift=-0.8mm]\x,0) -- ([yshift=0.8mm]\x,0);
  \foreach \x [count=\i from 0] in {0,10,...,60} \draw[ultra thick] ([yshift=-1.5mm]\x,0) coordinate (a\i) -- ([yshift=1.5mm]\x,0);
  \foreach \x [count=\i from 0] in {0,1,2,3,4,5,6} {
    \node[below=2mm of a\i,inner sep=0pt,font=\small] {$\num{\x}$};
  }

\end{tikzpicture}\\
Dans une unité, il y a $\ldots$ intervalles, donc le pas est de $\ldots$.\\
Le point $A$ a pour abscisse $\dfrac{7}{3}$, il faut donc avancer de $\ldots$ pas à partir de l'origine pour placer le point $A$.\\On effectue la même démarche pour les points $B$ et $C$. 


\item Placer le point $~B\left(\dfrac{32}{6}\right)$.


\bigskip
  \begin{tikzpicture}[x=2.5mm]
  \draw[-{Latex[round]},thick] (0,0) -- (61,0);
  \foreach \x in {0,1.66667,...,60} \draw[thick] ([yshift=-0.8mm]\x,0) -- ([yshift=0.8mm]\x,0);
  \foreach \x [count=\i from 0] in {0,10,...,60} \draw[ultra thick] ([yshift=-1.5mm]\x,0) coordinate (a\i) -- ([yshift=1.5mm]\x,0);
  \foreach \x [count=\i from 0] in {3,4,5,6,7,8,9} {
    \node[below=2mm of a\i,inner sep=0pt,font=\small] {$\num{\x}$};
  }

  
    
\end{tikzpicture} \\
Dans une unité, il y a $\ldots$ intervalles, donc le pas est de $\ldots$.\\Le point $B$ a pour abscisse $\dfrac{32}{6}$, il faut donc avancer de $\ldots$ pas à partir de l'origine pour placer le point $B$. \\
L'origine n'étant pas visible, pour aller à $3$ unités il faut avancer de $\ldots\ldots\ldots$ pas . \\Il nous reste donc à avancer de $\ldots\ldots\ldots$ à partir de $3$ pour placer le point $B$.
\end{enumerate}
\end{Exemple}

\vspace{0.8cm}
\subsection*{IV. Fractions égales}
\addcontentsline{toc}{subsection}{IV. Fractions égales}

\begin{Propriété}
Le nombre représenté par une fraction ne change pas lorsque l'on multiplie ou lorsque l'on divise le numérateur et le dénominateur de celle-ci par un même nombre non nul.
\end{Propriété}


\begin{Exemples}
\begin{itemize}[label=\textcolor{Couleur8}{$\bullet$}]
\item $\dfrac{2}{5}= \dfrac{\ldots\ldots}{\ldots\ldots} = \dfrac{\ldots}{20}$

\item $ \dfrac{18}{12} = \dfrac{\ldots\ldots}{\ldots\ldots} = \dfrac{3}{\ldots}$.\\
 Ici, on obtient une fraction \textbf{simplifiée} car le numérateur et le dénominateur sont plus petits que ceux de la fraction initiale.

\item On peut aussi simplifier une fraction avec des décompositions : $ \dfrac{21}{14} = \dfrac{\ldots\ldots}{\ldots\ldots} = \dfrac{\ldots}{2}$.
\end{itemize}
\end{Exemples}


\newpage

\subsection*{V. Comparaisons de fractions}
\addcontentsline{toc}{subsection}{V. Comparaisons de fractions}



\begin{Propriétés*}[c]
\leavevmode\vspace{-\baselineskip}
\begin{itemize}[label=\textcolor{Couleur1}{$\bullet$}]
\item Si deux fractions ont le même dénominateur, alors la fraction avec le plus grand numérateur est la fraction la plus grande.

\item Si deux fractions ont le même numérateur, alors la fraction avec le plus grand dénominateur est la fraction la plus petite.
\end{itemize}
\end{Propriétés*}

\begin{Exemples}
\begin{itemize}[label=\textcolor{Couleur8}{$\bullet$}]
\item Comparer $\dfrac{6}{7}$ et $\dfrac{2}{7}$. 

\ligne

\ligne

\ligne

\item Comparer $\dfrac{7}{12}$ et $\dfrac{7}{9}$. 

\ligne

\end{itemize}
\end{Exemples}

\begin{Propriétés*}[c]
\leavevmode\vspace{-\baselineskip}
\begin{itemize}[label=\textcolor{Couleur1}{$\bullet$}]
\item Une fraction est égale à 1 si et seulement si son numérateur est égal à son dénominateur.

\item Une fraction est strictement inférieure à 1 si et seulement si son numérateur est strictement inférieur à son dénominateur.

\item Une fraction est strictement supérieure à 1 si et seulement si son numérateur est strictement supérieur à son dénominateur.
\end{itemize}
\end{Propriétés*}

\begin{Exemples}
\vspace{-0.5cm}
\begin{multicols}{2}
\begin{itemize}[label=\textcolor{Couleur8}{$\bullet$}]
\item $\dfrac{6}{10} \ldots 1$ car \ligne

\item $\dfrac{9}{5} \ldots 1$ car \ligne
\end{itemize}
\end{multicols}
\begin{itemize}[label=\textcolor{Couleur8}{$\bullet$}]
\item $\dfrac{142}{142} \ldots 1$ car \ligne

\end{itemize}
\end{Exemples}

\begin{Exemple}
En utilisant la comparaison d'une fraction avec 1, on peut comparer deux fractions de dénominateurs différents dans certains cas. \\
Pour $\dfrac{5}{4}$ et $\dfrac{3}{7}$, on sait $\dfrac{5}{4} \ldots 1$ car $\ldots \ldots$. et que $\dfrac{3}{7} \ldots 1$. Donc \ligne

\end{Exemple}


\begin{Exemple}
\begin{enumerate}
\item Décomposer la fraction $\dfrac{9}{2}$ sous la forme d'un entier et d'une fraction inférieure à 1. 

\ligne

\item Donner un encadrement à l'unité de $\dfrac{9}{2}$.

\ligne

\end{enumerate}
\end{Exemple}
